\documentclass{article}
\usepackage{amsmath,hyperref}
\title{Door retirement needs either proven quiescence or passage identity}
\author{}
\date{}
\begin{document}
\maketitle
\paragraph{Question and source audit.}
The stock elapsed-time rule in the prior
\href{https://github.com/eternal-november/eos-research-programmes/pull/509}{retired-session replay}
can distinguish a delayed older stamp while a newer request is
logged. It does not identify the passage that generated a message.
The pinned planner passes a persistent requester ID, not a
per-passage identifier. The close source stamps the publication
with its current clock on each retry. The receiver sees only the
resulting named request, not the identity of the original phase.

A casually assumed two-thread race is not a source argument here.
The pinned default adapter creates a single event-loop worker and
passes it both to the transport and the fleet. Fleet robot context
and the plan's legacy shim use that same worker. The transport
starts a spin thread, but it schedules \texttt{spin\_some} onto
the worker rather than running the timer body on the spin thread.
The sequenced phase signals completion using the same worker;
the close observable terminates on completion or failure and
its \texttt{finally} resets the timer. Destruction of the legacy
shim's subscription guard explicitly unsubscribes while its
legacy instance is still held. A queued timer body attempts a
weak lock, so a genuinely released old phase cannot publish.
Replacing or stopping an execution is not itself a proof that
every external strong reference is gone. Operator kill, restored
or externally shared phases, alternative executors and delivery
after retirement were not exhaustively verified. In particular,
the replay does not assert that two old and new phases actually
overlap in the default task runtime.

\paragraph{Native-source controls under the explicit lifetime model.}
The replay executes the original close timer lambda, request
publication, local hold-clear helper, and supervisor methods.
One serial fixture retires the old phase and destroys its final
owner before a queued callback; that callback emits nothing.
In a separate retained-phase counterfactual, both the old retry
after the new hold and a legitimate current release produce the
same request (including the fresh timestamp). A second holder
prevents immediate CLOSE, masking the first removal. Once that
holder releases normally, the supervisor's final CLOSE reveals
the loss of the original live hold.

In the modeled generation-aware comparator, binding the producer
to the current passage prevents the old release, while the
legitimate same-timestamp release still succeeds. The data
needed for that binding are supplied by the fixture, not detected
by an unmodified node. The same test also runs the original
local name-only release helper; that helper clears a held marker
even if invoked for the earlier occurrence of the same name.
A scoped repair must qualify that call too. The idle heartbeat
cleanup is another producer that uses a stable requester ID
and a fresh stamp, so this is not a claim of comprehensive
repair of every source path. Source retirement also does
not settle ingress order: the supervisor accepts equality,
so a delayed old close with a stamp equal to the new open is
order-dependent in the fixture. Proving actual transport
ordering, strict time separation, or identity is separate.

\paragraph{Connection to the clearance consumer.}
The unchanged \href{https://github.com/eternal-november/eos-research-programmes/pull/809}{gate receipt checker}
can consume a retain-through-use fact from a suitable source.
The replay exercises only the logical hold factor of that
fact. It neither samples the motion of a door nor establishes
the physical aperture premise highlighted by
\href{https://github.com/eternal-november/eos-research-programmes/pull/810}{the retained-actuation study}.
This source boundary is valuable as a proof obligation
for adapters, not an observed stock incident, deployed guard,
or result closing I1/I2.
\end{document}
